\documentclass[11pt]{article}
\usepackage[margin=1.1in]{geometry}
\usepackage{amsmath,amssymb,amsthm,mathtools,booktabs,array,tabularx,longtable}
\usepackage[dvipsnames]{xcolor}
\usepackage[colorlinks=true,linkcolor=MidnightBlue,citecolor=MidnightBlue,urlcolor=MidnightBlue]{hyperref}
\usepackage{fancyhdr}
\pagestyle{fancy}\fancyhf{}\renewcommand{\headrulewidth}{0pt}
\fancyfoot[R]{\thepage}
\fancyfoot[L]{\scriptsize\color{gray}Generated by the Transport Lab from approved parts of its store; not a deposited paper. Papers of record: \href{https://zenodo.org/communities/transport-papers}{Zenodo, Transport Papers}}
\fancypagestyle{plain}{\fancyhf{}\fancyfoot[R]{\thepage}\fancyfoot[L]{\scriptsize\color{gray}Generated by the Transport Lab from approved parts of its store; not a deposited paper. Papers of record: \href{https://zenodo.org/communities/transport-papers}{Zenodo, Transport Papers}}\renewcommand{\headrulewidth}{0pt}}
\newtheorem{theorem}{Theorem}[section]
\newtheorem{proposition}[theorem]{Proposition}
\newtheorem{lemma}[theorem]{Lemma}
\newtheorem{corollary}[theorem]{Corollary}
\newtheorem{conjecture}[theorem]{Conjecture}
\newtheorem{refutation}[theorem]{Refutation}
\newtheorem{prediction}[theorem]{Prediction}
\newtheorem{principle}[theorem]{Principle}
\newtheorem{observation}[theorem]{Observation}
\newtheorem{axiom}[theorem]{Axiom}
\theoremstyle{definition}
\newtheorem{definition}[theorem]{Definition}
\newtheorem{fact}[theorem]{Fact}
\newtheorem{assembled}[theorem]{Assembly}
\newtheorem{postulate}[theorem]{Postulate}
\newtheorem{premise}[theorem]{Premise}
\newtheorem{hypothesis}[theorem]{Hypothesis}
\newtheorem{remark}[theorem]{Remark}
\newtheorem{example}[theorem]{Example}
\theoremstyle{plain}
\setlength{\marginparwidth}{0.8in}
\newcommand{\changedmark}[2]{\marginpar{\raggedright\scriptsize\color{#1}#2}}
\providecommand{\Aa}{}\renewcommand{\Aa}{\mathcal{A}}
\providecommand{\Ad}{}\renewcommand{\Ad}{\widetilde{\mathrm{Ad}}}
\providecommand{\Aev}{}\renewcommand{\Aev}{\mathcal{A}^{+}}
\providecommand{\Aod}{}\renewcommand{\Aod}{\mathcal{A}^{-}}
\providecommand{\B}{}\renewcommand{\B}{\mathrm{B}}
\providecommand{\C}{}\renewcommand{\C}{\mathbb{C}}
\providecommand{\CAR}{}\renewcommand{\CAR}{\mathrm{CAR}}
\providecommand{\CS}{}\renewcommand{\CS}{{\mathbb{S}^{+}}}
\providecommand{\CSm}{}\renewcommand{\CSm}{{\mathbb{S}^{-}}}
\providecommand{\Cl}{}\renewcommand{\Cl}{\mathrm{Cl}}
\providecommand{\Con}{}\renewcommand{\Con}{\mathcal{C}}
\providecommand{\Dl}{}\renewcommand{\Dl}{\Delta}
\providecommand{\EN}{}\renewcommand{\EN}{E_{N}}
\providecommand{\End}{}\renewcommand{\End}{\mathrm{End}}
\providecommand{\F}{}\renewcommand{\F}{\mathbb{F}}
\providecommand{\G}{}\renewcommand{\G}{\Gamma}
\providecommand{\HH}{}\renewcommand{\HH}{\mathbb{H}}
\providecommand{\IS}{}\renewcommand{\IS}{I\sigma_3}
\providecommand{\J}{}\renewcommand{\J}{\mathbf{J}}
\providecommand{\LH}{}\renewcommand{\LH}{L_{H}}
\providecommand{\Mat}{}\renewcommand{\Mat}{\mathrm{M}}
\providecommand{\Mo}{}\renewcommand{\Mo}{M^{\circ}}
\providecommand{\Mp}{}\renewcommand{\Mp}{M^{+}}
\providecommand{\Msun}{}\renewcommand{\Msun}{M_{\odot}}
\providecommand{\Mt}{}\renewcommand{\Mt}{\widetilde{M}}
\providecommand{\Ncone}{}\renewcommand{\Ncone}{\mathcal{N}}
\providecommand{\Neff}{}\renewcommand{\Neff}{N_{\mathrm{eff}}}
\providecommand{\PCS}{}\renewcommand{\PCS}{\mathbb{P}{\mathbb{S}^{+}}}
\providecommand{\PT}{}\renewcommand{\PT}{\mathbb{PT}}
\providecommand{\Pin}{}\renewcommand{\Pin}{\mathrm{Pin}}
\providecommand{\Q}{}\renewcommand{\Q}{\mathbb{Q}}
\providecommand{\R}{}\renewcommand{\R}{\mathbb{R}}
\providecommand{\Ric}{}\renewcommand{\Ric}{\mathrm{Ric}}
\providecommand{\SBH}{}\renewcommand{\SBH}{S_{\mathrm{BH}}}
\providecommand{\Sc}{}\renewcommand{\Sc}{\mathbb{S}}
\providecommand{\Sch}{}\renewcommand{\Sch}{\mathrm{P}}
\providecommand{\Sig}{}\renewcommand{\Sig}{\Sigma}
\providecommand{\Som}{}\renewcommand{\Som}{S_\omega}
\providecommand{\Sp}{}\renewcommand{\Sp}{\mathbb{S}}
\providecommand{\Sph}{}\renewcommand{\Sph}{\mathrm{S}}
\providecommand{\Spi}{}\renewcommand{\Spi}{S_\pi}
\providecommand{\Spin}{}\renewcommand{\Spin}{\mathrm{Spin}}
\providecommand{\Sth}{}\renewcommand{\Sth}{S_{\mathrm{th}}}
\providecommand{\T}{}\renewcommand{\T}{\mathbb{T}}
\providecommand{\TH}{}\renewcommand{\TH}{T_{\mathrm{H}}}
\providecommand{\Tr}{}\renewcommand{\Tr}{\operatorname{Tr}}
\providecommand{\UHF}{}\renewcommand{\UHF}{\mathrm{UHF}}
\providecommand{\Z}{}\renewcommand{\Z}{\mathbb{Z}}
\providecommand{\Zc}{}\renewcommand{\Zc}{\mathrm{Z}}
\providecommand{\cmg}{}\renewcommand{\cmg}{\mathrm{cm^{2}\,g^{-1}}}
\providecommand{\dS}{}\renewcommand{\dS}{\mathrm{dS}}
\providecommand{\dsc}{}\renewcommand{\dsc}{\mathrm{dsc}}
\providecommand{\dtPG}{}\renewcommand{\dtPG}{dt}
\providecommand{\dtS}{}\renewcommand{\dtS}{dt_{\mathrm{S}}}
\providecommand{\ellv}{}\renewcommand{\ellv}{\ell}
\providecommand{\er}{}\renewcommand{\er}{e_r}
\providecommand{\et}{}\renewcommand{\et}{e_t}
\providecommand{\ev}{}\renewcommand{\ev}{\mathrm{ev}}
\providecommand{\gl}{}\renewcommand{\gl}{\mathfrak{gl}}
\providecommand{\hatV}{}\renewcommand{\hatV}{\widehat{V}}
\providecommand{\hb}{}\renewcommand{\hb}{\bar{h}}
\providecommand{\hf}{}\renewcommand{\hf}{\hat{f}}
\providecommand{\id}{}\renewcommand{\id}{\mathrm{id}}
\providecommand{\kpc}{}\renewcommand{\kpc}{\mathrm{kpc}}
\providecommand{\lP}{}\renewcommand{\lP}{\ell_{P}}
\providecommand{\mD}{}\renewcommand{\mD}{\mathbf{m}_{D}}
\providecommand{\mb}{}\renewcommand{\mb}{\bar m}
\providecommand{\mnu}{}\renewcommand{\mnu}{\Sigma m_{\nu}}
\providecommand{\nInf}{}\renewcommand{\nInf}{n_\infty}
\providecommand{\nO}{}\renewcommand{\nO}{n_0}
\providecommand{\ninf}{}\renewcommand{\ninf}{n_{\infty}}
\providecommand{\nz}{}\renewcommand{\nz}{n_{0}}
\providecommand{\proj}{}\renewcommand{\proj}{\operatorname{proj}}
\providecommand{\pv}{}\renewcommand{\pv}{\pi}
\providecommand{\rad}{}\renewcommand{\rad}{\mathrm{rad}}
\providecommand{\sgn}{}\renewcommand{\sgn}{\operatorname{sgn}}
\providecommand{\shA}{}\renewcommand{\shA}{\mathcal{A}}
\providecommand{\shC}{}\renewcommand{\shC}{\mathcal{C}}
\providecommand{\shE}{}\renewcommand{\shE}{\mathcal{E}}
\providecommand{\shJ}{}\renewcommand{\shJ}{\mathcal{J}}
\providecommand{\shK}{}\renewcommand{\shK}{\mathcal{K}}
\providecommand{\shL}{}\renewcommand{\shL}{L}
\providecommand{\shW}{}\renewcommand{\shW}{\mathcal{W}}
\providecommand{\sll}{}\renewcommand{\sll}{\mathfrak{sl}}
\providecommand{\so}{}\renewcommand{\so}{\mathfrak{so}}
\providecommand{\su}{}\renewcommand{\su}{\mathfrak{su}}
\providecommand{\supp}{}\renewcommand{\supp}{\operatorname{supp}}
\providecommand{\tr}{}\renewcommand{\tr}{\operatorname{tr}}
\providecommand{\uu}{}\renewcommand{\uu}{\mathfrak{u}}
\providecommand{\br}{}\renewcommand{\br}[1]{\langle #1\rangle}
\providecommand{\bra}{}\renewcommand{\bra}[1]{\langle #1|}
\providecommand{\file}{}\renewcommand{\file}[1]{\texttt{#1}}
\providecommand{\hp}{}\renewcommand{\hp}[2]{( #1 \,|\, #2 )}
\providecommand{\ip}{}\renewcommand{\ip}[2]{\langle #1, #2\rangle}
\providecommand{\ket}{}\renewcommand{\ket}[1]{|#1\rangle}
\providecommand{\rel}{}\renewcommand{\rel}[2]{D\bigl(#1 \,\|\, #2\bigr)}
\providecommand{\rev}{}\renewcommand{\rev}[1]{\widetilde{#1}}
\providecommand{\slashed}{}\renewcommand{\slashed}[1]{{#1}\mkern-10mu/}
\title{\bfseries The origin of mass}
\author{Leon Butler\\ \normalsize\href{mailto:lbutler.math@gmail.com}{\texttt{lbutler.math@gmail.com}}\\ \normalsize ORCID: \href{https://orcid.org/0009-0006-4449-5617}{0009-0006-4449-5617}}
\date{Generated by the Transport Lab, 2 November 2026}
\begin{document}
\maketitle
\begin{abstract}
\normalsize
Where does mass come from in this framework? The completion has one built-in scale, the frozen point set by the cosmological constant, and mass is read as how strongly a particle grips it.

Two massless conformal spinors (twistors) make a massive composite whose mass measures how far apart their light-like directions point, and the trail shows the two readings agree exactly. A second reading turns mass into one number a field carries, its conformal weight.

The trail then meets the subject's main refutation: tying every mass to the cosmological scale restates the hierarchy rather than explaining it. It closes on a clear record of which routes to the absolute mass scale fail.
\end{abstract}
\newpage
\noindent \textbf{Authorship and disclosure. }The founding idea of the series is the author's: that division by zero in a Clifford algebra should be read not as a prohibition but as a transport into a higher-dimensional algebra, and that the transported division need not be accompanied by an inverse operation. The direction of investigation and the critical review of the results are likewise the author's. The mathematical development and the drafting of the text were carried out in extended collaboration with Claude Fable 5 (Anthropic), an AI system.

\noindent\fbox{\parbox{\dimexpr\linewidth-2\fboxsep-2\fboxrule}{\textbf{This is not a deposited paper.} It was generated by the Transport Lab on 2 November 2026 and assembled from approved parts of its store. It follows results across several papers of the series; each is stated with its proof and credited to the paper that first proved it. The official versions of those papers are on Zenodo, in the Transport Papers community, \url{https://zenodo.org/communities/transport-papers}.}}
\newpage

\section*{Introduction}
The walk first fixes the one scale the framework has, then builds mass from conformal spinors and from weights, joins the two, and ends with what the mechanism does not explain.

Results a step rests on but that are not on the trail are listed at the end.

\section{Dark energy and the frozen point}

Mass needs a scale to be measured against, so the walk begins with the frozen point of the dark-energy subject.

On flat spacetime, de Sitter space is the gravitational field $\hb = \sigma$, a pure dilation by the function $\sigma = 1 - \Lambda Q(x)/12$.

Its frozen point is the unique constant vector \[ I\_\Lambda = n\_\infty + \frac{\Lambda}{6}\,n\_0 \] of squared norm $\Lambda/3$, with $n_0$ and $n_\infty$ the two null directions the completion adds. With the cosmological constant switched off it is the point at infinity $n_\infty$.

\begin{lemma}[The de Sitter scale and its frozen point]\label{c:mass.de-sitter-scale-frozen-point}
De Sitter space, the universe with only a positive cosmological constant, needs no curved background: on ordinary flat spacetime it is a gravitational field that rescales every length by one simple function of position. This lemma finds the frozen point of that field and shows that it is a single fixed vector of the larger conformal space. When the cosmological constant is switched off, that vector becomes the familiar point at infinity of flat space's conformal model.
\emph{First proved in} T11, Lem 5.1 (\href{https://doi.org/10.5281/zenodo.22722565}{doi:10.5281/zenodo.22722565}).
\end{lemma}
\noindent\emph{Rests on:} \ref{c:mass.imported-flat-model-twistor-chart-infinity}, \ref{c:mass.imported-frozen-point-tractor-connection}, \ref{c:mass.imported-thomas-operator}.

\section{The origin of mass}

With the scale in hand, these results build mass from conformal spinors and from conformal weights, and test what the construction explains.

Before using the frozen point, recall how two conformal spinors make mass. A single conformal spinor carries null momentum and is massless.

Two of them with spinor parts $\pi^{1}$ and $\pi^{2}$ have \[ m\^{}{2} = 2\,\bigl|\br{\pi\^{}{1}\pi\^{}{2}}\bigr|\^{}{2}, \] the squared symplectic area between the two spinor parts. Aligned spinors give a massless composite, and mass grows with the misalignment.

\begin{proposition}[Mass as symplectic area: the two-spinor formula]\label{c:mass.mass-symplectic-area-two-twistor-formula}
A single spinor of the conformal geometry describes a massless particle travelling at light speed. Two of them together can describe a massive particle, and its mass depends only on how far apart their two light-like directions point. When they line up the composite stays massless; the more they disagree, the heavier it is.
\emph{First proved in} T11, Prop 4.1 (\href{https://doi.org/10.5281/zenodo.22722565}{doi:10.5281/zenodo.22722565}).
\end{proposition}

To read mass off the frozen point we need to know which symmetries it breaks.

For a non-null $I$ the stabiliser in $\mathfrak{sl}(\CS)$ is $10$-dimensional, and the five generators $A(e_{a}\wedge I)$ span its $5$-dimensional complement. For the frozen point with $\Lambda > 0$ these five are the de Sitter translations.

\begin{lemma}[The five broken generators]\label{c:mass.five-broken-generators}
The cosmological scale picks out a direction that breaks the conformal symmetry of spacetime down to de Sitter symmetry. Of the fifteen conformal symmetries ten survive and five are broken. This result identifies the five broken ones exactly: they are the de Sitter translations, the symmetries whose charges make up a composite's momentum.
\emph{First proved in} T11, Lem 5.8 (\href{https://doi.org/10.5281/zenodo.22722565}{doi:10.5281/zenodo.22722565}).
\end{lemma}
\noindent\emph{Rests on:} \ref{c:mass.de-sitter-scale-frozen-point}, \ref{c:mass.imported-flat-model-twistor-chart-infinity}, \ref{c:mass.imported-frozen-point-tractor-connection}.

There are now two readings of mass: the grip on the frozen point and the misalignment of two conformal spinors.

This result joins them exactly: \[ m\^{}{2} = 2\,|B\_{\Lambda}|\^{}{2} + \frac{\Lambda}{3}\,C, \] where $B_{\Lambda} = \br{\pi^{1}\pi^{2}} + \tfrac{\Lambda}{6}[\omega^{1}\omega^{2}]$ is the deformed bracket and $C$ depends on how the constituents overlap. At $\Lambda = 0$ the two-spinor formula returns.

\begin{theorem}[The exact join]\label{c:mass.exact-join}
The framework reads mass in two ways: from how a composite grips the one built-in cosmological scale, and from how misaligned its two constituents are. This theorem joins the two readings exactly. The mass squared is twice the squared size of a slightly deformed misalignment, plus a tiny correction, proportional to the cosmological constant, that depends on how the constituents overlap. Without a cosmological constant the familiar two-constituent area formula returns.
\emph{First proved in} T11, Thm 5.6 (\href{https://doi.org/10.5281/zenodo.22722565}{doi:10.5281/zenodo.22722565}).
\end{theorem}
\noindent\emph{Rests on:} \ref{c:correspondence.incidence-relation-chart}, \ref{c:mass.chart-form-spin-isomorphism}, \ref{c:mass.de-sitter-scale-frozen-point}, \ref{c:mass.imported-flat-model-twistor-chart-infinity}, \ref{c:mass.mass-symplectic-area-two-twistor-formula}.

A different route to mass goes through fields rather than particles.

A density $\sigma$ of conformal weight $w$ coupled to the frozen point by $I \cdot D\sigma = 0$ obeys the Klein-Gordon equation, in four dimensions with \[ m\^{}{2} = -\,w(w+3)\,H\^{}{2} \] Every mass is then $\sqrt{\Lambda}$ times a factor fixed by the weight.

\begin{theorem}[The mass-weight dictionary]\label{c:mass.mass-weight-dictionary}
The framework has exactly one built-in scale, set by the cosmological constant. This theorem shows that a field coupled to that scale in the natural way behaves as a particle with a definite mass, fixed by a single number the field carries, its conformal weight. In that sense every mass is the cosmological scale multiplied by a number determined by the weight.
\emph{First proved in} T11, Thm 3.1 (\href{https://doi.org/10.5281/zenodo.22722565}{doi:10.5281/zenodo.22722565}).
\end{theorem}
\noindent\emph{Rests on:} \ref{c:mass.imported-frozen-point-tractor-connection}, \ref{c:mass.imported-thomas-operator}.

The next step maps out which masses this rule allows.

Real weights give $0 \le m \le 3H/2$, exactly for $w$ in $[-3, 0]$. Anything heavier needs the complex weight $w = -\tfrac32 \pm i\nu$, the principal series, and the two regimes match the two standard series of de Sitter representations.

\begin{corollary}[The weight polynomial and the two series]\label{c:mass.weight-polynomial-two-series}
The framework's mass rule turns one number, the weight, into a mass. This result maps out what the rule allows. Real weights only give masses up to a tiny threshold set by the expansion rate of the universe; anything heavier needs a complex weight, with a fixed real part and an imaginary part that carries the mass. These two regimes match the two standard families of particles in de Sitter space.
\emph{First proved in} T11, Cor 3.2 (\href{https://doi.org/10.5281/zenodo.22722565}{doi:10.5281/zenodo.22722565}).
\end{corollary}
\noindent\emph{Rests on:} \ref{c:mass.imported-unitary-representations-de-sitter-group}.

Combining the weight rule with the exact join assigns every two-spinor composite a weight.

For $\Lambda > 0$ a composite heavier than $3H/2$ has \[ w = -\frac32 \pm i\nu, \] the same real part for every such particle. Its mass lies in the imaginary part, which at leading order is $\sqrt{2}\,|B_{\Lambda}|/H$, an enormous number for every known particle.

\begin{proposition}[The weight of a composite]\label{c:mass.weight-composite}
With a positive cosmological constant, each composite particle can be given a weight: the number that would produce its mass under the framework's mass rule. For anything heavier than a tiny threshold set by the expansion of the universe, that weight is complex. Its real part is the same for every such particle, and its imaginary part is essentially the mass measured in cosmological units, an enormous number for every known particle.
\emph{First proved in} T11, Prop 5.9 (\href{https://doi.org/10.5281/zenodo.22722565}{doi:10.5281/zenodo.22722565}).
\end{proposition}
\noindent\emph{Rests on:} \ref{c:mass.exact-join}, \ref{c:mass.weight-polynomial-two-series}.

The size of a mass is one question; which field carries the mass term is another.

Under the two-spinor $U(2)$ the mass bracket transforms by the determinant, \[ \br{\pi\^{}{1}\pi\^{}{2}} \mapsto \det(U)\,\br{\pi\^{}{1}\pi\^{}{2}}, \] so it is an $SU(2)$ singlet of hypercharge $+1$. It carries the quantum numbers of $e^{c}$, not of a weak doublet such as the Higgs.

\begin{proposition}[The mass bracket is a charged singlet]\label{c:standard-model.mass-bracket-charged-singlet}
For two conformal spinors, the series' mass is the size of a single complex number, the mass bracket. This result works out how that number changes under the internal symmetries of a system of two conformal spinors. It is unchanged by the $SU(2)$ part and only picks up a phase from the $U(1)$ part, so, with the paper's global form and hypercharge, it carries the charges of the conjugate electron field $e^{c}$, a weak singlet of hypercharge $+1$, rather than those of a weak doublet such as the Higgs.
\emph{First proved in} T12, Prop 3.2 (\href{https://doi.org/10.5281/zenodo.22722669}{doi:10.5281/zenodo.22722669}).
\end{proposition}
\noindent\emph{Rests on:} \ref{c:mass.mass-symplectic-area-two-twistor-formula}, \ref{c:standard-model.global-form-quotient-determined}, \ref{c:standard-model.imported-pair-systems-two-chirality-halves}.

It was natural to hope that tying every mass to the one cosmological scale would explain why particles weigh what they do.

The programme records that hope as refuted: coupling to the frozen point restates the hierarchy as the size of the imaginary weights in Hubble units. The old puzzle returns in a new form.

\begin{refutation}[The mass dictionary, coupling to the frozen point, explains the observed masses]\label{c:refuted.the-mass-dictionary-coupling-to-the-frozen-point}
It was natural to hope that tying every field to the single cosmological scale would explain why particles have the masses they do. It does not: that scale is so small that ordinary particles need enormous weights, so the old puzzle of why matter is so heavy compared with it simply comes back in a new form.
\end{refutation}

The remaining question is whether anything else in the framework fixes the size of those weights.

Every structural route tried, including the pion-scale $N^{2/3}$ scaling, fails, so the absolute mass hierarchy stays one unexplained number. It is recorded as a clean negative result, at the level of a conjecture rather than a theorem.

\begin{conjecture}[The absolute scale: a clean negative result]\label{c:families.absolute-scale-clean-negative-result}
In this framework a particle's mass is coded by a number, its weight, measured against the tiny scale set by the cosmological constant. For ordinary particles that number is enormous, around $10^{38}$ for the electron, and nothing in the framework yet explains why. Several natural attempts were tried and all fail, one of them because it only rediscovers a coincidence first noticed in 1967. What remains is a clear record of which routes do not work.
\emph{First proved in} T13, Obs 7.4 (\href{https://doi.org/10.5281/zenodo.22722845}{doi:10.5281/zenodo.22722845}).
\end{conjecture}

\section{Results used from elsewhere}
The proofs above use the following results of the series, stated briefly; each is proved in the paper named.

\begin{fact}[The incidence relation is the chart]\label{c:correspondence.incidence-relation-chart}
The Pl\"ucker bivector of the incidence line of a map $A$ is exactly the completion's null point embedding $X(m(A))$, so Penrose's incidence relation is the completion's affine chart, with a determinant as its quadratic term.
\emph{Source:} T9, Thm 4.2 (\href{https://doi.org/10.5281/zenodo.22722290}{doi:10.5281/zenodo.22722290}).
\end{fact}

\begin{fact}[The chart form of the spin isomorphism]\label{c:mass.chart-form-spin-isomorphism}
In the conformal spinor chart, the dilation, the translations and the special conformal generators correspond to the grading matrix and to the upper and lower off-diagonal blocks $b_x$ and $c_x$.
\emph{Source:} T11, Lem 5.3 (\href{https://doi.org/10.5281/zenodo.22722565}{doi:10.5281/zenodo.22722565}).
\end{fact}

\begin{fact}[Imported: the flat model, the conformal spinor chart, and the infinity bivector]\label{c:mass.imported-flat-model-twistor-chart-infinity}
Imported: on the flat model the point embedding is null, $V_\C = \Lambda^2\CS$ with null meaning simple, $\mathfrak{sl}(\CS) \cong \so(\Lambda^2\CS, \varepsilon)$, and for $\Lambda = 0$ the infinity bivector is constant and simple.
\emph{Source:} T11, Fact 2.3 (\href{https://doi.org/10.5281/zenodo.22722565}{doi:10.5281/zenodo.22722565}).
\end{fact}

\begin{fact}[Imported: the frozen point and the tractor connection]\label{c:mass.imported-frozen-point-tractor-connection}
For a vacuum Einstein field on flat spacetime the tractor gauge field has a parallel frozen point $I = \nInf + \tfrac{\Lambda}{6}\nO$ of squared norm $\Lambda/3$, with stabiliser $SO(1,4)$ for $\Lambda > 0$.
\emph{Source:} T11, Fact 2.1 (\href{https://doi.org/10.5281/zenodo.22722565}{doi:10.5281/zenodo.22722565}).
\end{fact}

\begin{fact}[Imported: the Thomas operator]\label{c:mass.imported-thomas-operator}
The Thomas operator sends a weight-$w$ density to a conformally invariant tractor, and at $w = 1$ it produces the scale tractor, which is the frozen point for an Einstein scale.
\emph{Source:} T. Y. Thomas, On conformal geometry, Proc. Natl. Acad. Sci. USA 12 (1926), 352--359.; T. N. Bailey, M. G. Eastwood, and A. R. Gover, Thomas's structure bundle for conformal, projective and related structures, Rocky Mountain J. Math. 24 (1994), 1191--1217.; A. R. Gover, Almost Einstein and Poincar'e-Einstein manifolds in Riemannian signature, J. Geom. Phys. 60 (2010), 182--204..
\end{fact}

\begin{fact}[Imported: the unitary representations of the de Sitter group]\label{c:mass.imported-unitary-representations-de-sitter-group}
The scalar unitary representations of $SO(1,4)$, with $m^{2} = \Delta(3-\Delta)H^{2}$, comprise the principal series ($m^{2} \ge \tfrac94 H^{2}$), the complementary series ($0 < m^{2} < \tfrac94 H^{2}$) and the massless point, none with $m^{2} < 0$.
\emph{Source:} J. Dixmier, Repr'esentations int'egrables du groupe de De Sitter, Bull. Soc. Math. France 89 (1961), 9--41., \S\S3--4; A. Higuchi, Symmetric tensor spherical harmonics on the $N$-sphere and their application to the de Sitter group $SO(N,1)$, J. Math. Phys. 28 (1987), 1553--1566..
\end{fact}

\begin{fact}[The global form: the quotient determined]\label{c:standard-model.global-form-quotient-determined}
The global gauge group of the minimal massive composite is $S(U(2) \times U(3)) \cong (SU(2) \times SU(3) \times U(1))/\mathbb{Z}_{6}$, the smallest of the Standard Model's global forms.
\emph{Source:} T12, Thm 4.2 (\href{https://doi.org/10.5281/zenodo.22722669}{doi:10.5281/zenodo.22722669}).
\end{fact}

\begin{fact}[Imported: the pair systems and the two chirality halves]\label{c:standard-model.imported-pair-systems-two-chirality-halves}
Imported: the Fock module splits into the four-dimensional halves $\CS$ and $\CSm$, and every nonzero conformal spinor is a pure spinor, the Fock vacuum of its own three ladder pairs.
\emph{Source:} T12, Fact 2.2 (\href{https://doi.org/10.5281/zenodo.22722669}{doi:10.5281/zenodo.22722669}).
\end{fact}

\begin{prediction}[No fractionally charged colour singlet]\label{c:pred.z6-charge-lattice}
The $\mathbb{Z}_{6}$ charge lattice of the derived gauge group forbids any fractionally charged colour singlet.
\emph{Kill condition:} One confirmed fractionally charged colour-singlet particle.
\end{prediction}

\section*{Conclusion}
The trail has shown two readings of mass, the grip on the frozen point and the misalignment of two conformal spinors, joined exactly, with the conformal weight as a third language for the same quantity.

No postulate is needed on this stretch, and the charge assignment of the mass bracket rests on the series' derived global form and hypercharge. What the mechanism does not do is explain the observed masses: it restates the hierarchy as the size of the weights, which remains open.

\end{document}